%% ma: L10-B13-0049
%% lop: 10
%% chuong: 5
%% bai: 13
%% dang: 10.13.8
%% loai: DS
%% muc: [TH, NB, TH, TH]
%% dap_an: "ĐSĐS"
%% nguon: "chuyen-de-cac-so-dac-trung-cua-mau-so-lieu-khong-ghep-nhom-toan-10-chuong-trinh-moi.pdf (GV. Phan Nhật Linh) (P02), Bài 02 Dạng 1 Phần II Câu 4"
%% kien_thuc: "Lớp 10 Bài 13 (số trung bình, trung vị, ý nghĩa)"
%% trang_thai: chinh-thuc
%% ghi_chu: "Câu 6 gốc trùng hoàn toàn câu này, đã bỏ. Tự duyệt theo yêu cầu giáo viên 22/10/2026."
\begin{ex}
Bảng số liệu sau cho biết mức lương hằng năm của các cán bộ và nhân viên trong một công ty (đơn vị: nghìn đồng):
\begin{center}
\begin{tabular}{|c|c|c|c|}\hline
$20910$ & $76000$ & $20350$ & $20060$\\\hline
$21410$ & $20110$ & $21410$ & $21360$\\\hline
$20350$ & $21130$ & $20960$ & $125000$\\\hline
\end{tabular}
\end{center}
Xét tính đúng sai của các khẳng định sau:
\choiceTF{\True Mức lương trung bình là $\overline x=34\,087{,}5$ (nghìn đồng)}{Mức lương lớn nhất là $76\,000$}{\True Trung vị là $21\,045$ (nghìn đồng)}{Nên lấy mức lương trung bình làm giá trị đại diện cho mẫu số liệu}
\loigiai{a) Đúng: $\overline x=\dfrac{409\,050}{12}=34\,087{,}5$.

b) Sai: lớn nhất là $125\,000$.

c) Đúng: sắp xếp, giá trị thứ $6$ và $7$ là $20\,960$ và $21\,130$, nên $M_e=21\,045$.

d) Sai: mẫu có giá trị bất thường ($76\,000$; $125\,000$) làm số trung bình lớn hơn hẳn đa số; nên dùng trung vị làm đại diện.}
\end{ex}
